\chapter{Tools, contracts, and protocol boundaries}
\section{A tool is an interface to an operation}
A tool connects a model proposal to executable behavior. The behavior may be pure computation, a read from a service, or a change to an external system. Those categories have different failure and authorization requirements. A calculator can return a result without changing the environment. An inventory lookup reads changing state. A reservation call creates a lasting effect.

A good tool interface exposes the smallest operation the application needs. A function such as \texttt{reserve\_item(item\_id, date)} is easier to constrain than “run arbitrary database code.” Narrow interfaces make allowed behavior explicit, reduce the argument space, and give tests something concrete to exercise. They do not eliminate the need for checks inside the service.

Toolformer studied training a language model to select and use API calls \cite{toolformer}. This course instead focuses on application engineering around an already available model. Learning to propose useful calls and having permission to execute them are distinct concerns. A highly accurate model still needs a runtime with a clear authority boundary.

\section{Designing a contract}
A tool contract includes a name, purpose, input schema, output schema, side-effect description, error semantics, and resource limits. The description should distinguish similar tools. \texttt{lookup\_item} returns catalog information; \texttt{check\_availability} returns availability for a particular interval; \texttt{prepare\_reservation} produces a proposal; \texttt{commit\_reservation} changes the booking system. Hiding these differences behind one broad “equipment” tool makes correct use harder to verify.

Input validation should reject missing fields, unexpected fields where appropriate, wrong types, invalid ranges, and malformed identifiers. Domain validation then checks whether the authenticated user may act on the requested object. Never let a model-supplied \texttt{user\_id} override the authenticated identity without a separately authorized delegation mechanism.

Output contracts matter equally. A lookup should distinguish “item not found” from “service unavailable.” A timeout does not prove that the item does not exist. Include source identity and observation time when later decisions depend on freshness. Bound result size so that a huge response cannot overwhelm the next model context or the local process.

\section{Worked example: a small numeric tool}
Suppose the service needs a tool that adds two finite measurements within a teaching range. The following complete Python example illustrates type and magnitude checks. It is deliberately narrower than a general calculator.

\begin{Verbatim}[breaklines=true,breakanywhere=true,fontsize=\small,frame=single]
import math

LIMIT = 1_000_000.0

def bounded_add(a, b):
    for value in (a, b):
        if type(value) not in (int, float):
            raise TypeError("number required")
        if abs(value) > LIMIT:
            raise ValueError("operand outside range")
        if not math.isfinite(value):
            raise ValueError("finite operand required")
    result = a + b
    if not math.isfinite(result) or abs(result) > LIMIT:
        raise ValueError("result outside range")
    return result

assert bounded_add(2, 3) == 5
assert bounded_add(-LIMIT, LIMIT) == 0
\end{Verbatim}

The exact-type test excludes Boolean values, which Python otherwise treats as a subclass of integers. The range check precedes conversion-dependent floating-point checks, so an extremely large integer is rejected without trying to convert it to a float. The result is checked as well as the operands: two permitted operands can sum to an out-of-range result.

This function bounds arithmetic on values it receives. It does not protect an HTTP server against an enormous request body before parsing. Layered resource controls begin at input size, then apply to parsed structure, permitted operations, execution time, and output size. The repaired calculator foundation extends this idea to a restricted expression tree with depth and operation limits.

\section{Dispatch and error handling}
A dispatcher receives a parsed call and looks up the implementation in an explicit registry. It should not evaluate a model-supplied expression as arbitrary Python or dynamically import a function named by the model. An unknown tool name is a predictable validation result. Recording it makes the attempted behavior visible without executing it.

A tool exception needs classification. A transient read timeout might be retried within a budget. An invalid date should be repaired before retry. A permission denial should not trigger attempts to find another route to the same prohibited action. The runtime's error policy can return safe, bounded information to the model while recording more diagnostic detail in an appropriately protected log.

For effectful calls, retries are especially important. If the reservation service committed a booking but its response was lost, repeating the call can create a second booking unless the service supports deduplication. The request needs a stable operation identifier and a binding to the original payload. Chapters 7 and 13 develop this problem in detail.

\section{Gemini, the Agents SDK, and responsibility}
The Gemini native API and its OpenAI-compatible endpoint are different integration paths. In the course, the native path exposes the tool round-trip for inspection; the Agents SDK path uses an OpenAI-compatible Chat Completions model adapter. Google's compatibility documentation and the SDK model documentation describe the relevant interfaces \cite{gemini-openai,agents-models}. Compatibility should be tested for the exact features used; it is not an assertion that all OpenAI-hosted capabilities exist at another provider.

The SDK can manage model requests and local tool execution, but the tool body still owns domain validation. A framework turn limit is useful but does not replace all tool, time, and spending limits. Tracing configuration also deserves attention because traces can contain user input and tool output. The course disables provider tracing in its offline exercises and keeps live inference explicitly enabled by the instructor.

\section{What MCP standardizes}
The Model Context Protocol provides a shared protocol between hosts, clients, and servers for capabilities such as tools, resources, and prompts. In a typical arrangement, the host runs the application, a client maintains a connection, and a server exposes capabilities. The course uses the dated 2025-11-25 specification as its baseline \cite{mcp-spec}. These role names describe protocol responsibilities rather than separate physical machines in every deployment.

A tool is callable behavior; a resource supplies content; a prompt supplies a reusable interaction template. Distinguishing them helps the application decide how content enters context and how operations are invoked. A common protocol reduces bespoke integration work, but it does not make every server trustworthy or every returned document authoritative.

MCP roots are particularly easy to misunderstand. The versioned roots specification describes filesystem context advertised by a client; roots are not a security sandbox \cite{mcp-roots}. If a server process can read arbitrary files under its operating-system credentials, advertising one directory does not revoke those permissions. Actual access controls must exist in the server, operating system, or execution environment.

\section{A boundary audit}
Imagine a server advertises a policy-search tool. The tool returns a document containing a request to export the student's profile. The search result is data from an external source. It cannot enlarge the caller's authority. The application must keep the user's authorized task and authenticated identity outside the control of retrieved text.

Audit the interaction as a sequence: who selected the server, who authenticated it, which credentials it uses, what operation was requested, what arguments were accepted, what data returned, and which subsequent effect was permitted. A protocol can make this sequence interoperable without answering every security question in it.

\section{Exercises}
\begin{enumerate}
\item Specify separate contracts for availability lookup and reservation commitment. Include errors that must not be conflated.
\item Test \texttt{bounded\_add} with Boolean values, a huge integer, infinity, NaN, two large permitted operands, and a nonnumeric string. Explain the expected outcome for each.
\item Explain why a valid JSON object can still contain an unauthorized action. Identify where authorization obtains the trusted user identity.
\item A server advertises one filesystem root. What evidence would demonstrate that it cannot read outside that directory? Distinguish protocol metadata from an enforcement test.
\item Design a retry policy for a read timeout and for a commit timeout. Explain why the policies differ.
\end{enumerate}

\section{Further study and laboratory connection}
Lab 5 tests argument validation and the native call cycle. Lab 7 exercises the Agents SDK through the Gemini compatibility path. Read the exact protocol references \cite{mcp-spec,mcp-roots} before making claims about MCP permissions. Record the installed versions and features actually tested rather than relying on framework names as guarantees.

